\documentclass{article}
\usepackage[margin=1in]{geometry}
\usepackage{float}
\usepackage{listings}
\usepackage{xcolor}
\usepackage{graphicx}

% Adding the style definition back in to make sure your code blocks look great!
\definecolor{codegreen}{rgb}{0,0.6,0}
\definecolor{codegray}{rgb}{0.5,0.5,0.5}
\definecolor{codepurple}{rgb}{0.58,0,0.82}
\definecolor{backcolour}{rgb}{0.95,0.95,0.92}

\lstdefinestyle{mystyle}{
    backgroundcolor=\color{backcolour},   
    commentstyle=\color{codegreen},
    keywordstyle=\color{magenta},
    numberstyle=\tiny\color{codegray},
    stringstyle=\color{codepurple},
    basicstyle=\ttfamily\footnotesize,
    breakatwhitespace=false,         
    breaklines=true,                 
    captionpos=b,                    
    keepspaces=true,                 
    numbers=left,                    
    numbersep=5pt,                  
    showspaces=false,                
    showstringspaces=false,
    showtabs=false,                  
    tabsize=2
}
\lstset{style=mystyle}

\title{
    \vspace{-2cm}
    \Large Data Engineering 424  \\
    \Large Bayes Networks  \\
}
\author{
    \textbf{Name:} Jacob Johannes Mostert \\ 
    \textbf{Student Number:} 27038122
}
\date{21 February 2026}
\begin{document}
\maketitle

\section*{Declaration}
I declare that this report and the work presented in it are my own work. If I received any help or information from any source that contributed to the work or report, it has been properly acknowledged.

\section{Question 1.1: The Monty-Hall Problem}
The Monty-Hall problem consists of a TV game show participant choosing between three doors, aiming to find an expensive car rather than a smelly goat. This is modeled using three random variables: the initial choice ($I$), the car's actual location ($C$), and the door Monty opens ($M$). Because Monty's choice depends on both the car's location and the participant's initial choice, the appropriate Bayes Net structure is a collider: $I \rightarrow M \leftarrow C$.

\subsection*{(a) Joint Distribution}
Using the \texttt{emdw} library, the initial factors $P(I)$, $P(C)$, and $P(M|I,C)$ were defined. The joint distribution $P(M,I,C)$ was calculated using factor operations. 

\begin{lstlisting}[language=C++, caption=Constructing the joint distribution]
rcptr<Factor> jointICM = pI->absorb(pC)->absorb(pM_IC);
\end{lstlisting}

When printed, the resulting joint distribution table $P(I,C,M)$ contains 12 non-zero entries. The probabilities evaluated to either $\approx 0.0556$ or $\approx 0.1111$. The table proves that Monty's behavior changes depending on whether your initial choice was right or wrong.

\subsection*{(b) Solving the Monty-Hall Problem}
To determine if the player should switch doors, we calculate the distribution over the location of the car given specific evidence. Assuming the player initially chooses door 0 ($I=0$) and Monty opens door 1 ($M=1$), we observe and reduce these variables, then marginalize over $C$.

\begin{lstlisting}[language=C++, caption=Calculating car location probabilities]
rcptr<Factor> C_IM = jointICM->observeAndReduce({I, M}, {0, 1})->marginalize({C})->normalize();
\end{lstlisting}

Based on the output, the resulting probability distribution for the car's location is $P(C=0) \approx 0.333$ and $P(C=2) \approx 0.667$. The probability of the car being behind the initially chosen door (0) remains at approximately $33.3\%$, while the probability of it being behind the remaining closed door (2) increases to $66.7\%$. So, it makes sense for the player to switch to the other door.

\section{Question 1.2: Flow of Influence via Collider Nodes}
We investigate the conditional dependence and independence between random variables (RVs), specifically observing the flow of influence through the collider node at $M$.

\subsection*{(e) Flow of Influence with an Unreliable Reporter}
We introduce a new node $R$, which represents an unreliable report of which door Monty opened. $R$ is influenced directly by $M$, adding the edge $M \rightarrow R$. The factor $P(R|M)$ gives an $80\%$ chance of reporting the correct door and a $10\%$ chance each of erroneously reporting the other two doors.

We want to know if influence flows between $I$ and $C$ when $R$ is either observed or unobserved, keeping $M$ completely unobserved.

\begin{itemize}
    \item \textbf{When $R$ is unobserved:} Knowledge of $I$ tells us nothing about $C$. Our output confirms that $P(C|I=0)$ evaluates to $0.333$ for all three doors, which is identical to the prior $P(C)$. The marginal belief over $C$ remains unchanged. This is consistent with d-separation theory; the path $I \rightarrow M \leftarrow C$ is blocked at the collider $M$ because neither $M$ nor any of its descendants (like $R$) are observed.

    \item \textbf{When $R$ is observed:} If we observe $R$ (e.g., $R=1$), it provides evidence about $M$. Because $R$ is a descendant of the collider $M$, observing it opens the flow of influence between $I$ and $C$. The output demonstrates this dependence: observing just $R=1$ shifts the distribution to $P(C=0)=0.45$, $P(C=1)=0.10$, and $P(C=2)=0.45$. If we then set $I=0$, the probabilities changed to $P(C=0) \approx 0.333$, $P(C=1) \approx 0.074$, and $P(C=2) \approx 0.593$. They are no longer independent.
\end{itemize}

\begin{lstlisting}[language=C++, caption=Testing influence with unobserved and observed R]
// Observe I = 0, R unobserved
rcptr<Factor> margC_given_I2 = jointICMR->observeAndReduce({I}, {0})->marginalize({C})->normalize();

// Observe R = 1, I unobserved
rcptr<Factor> margC_given_R = jointICMR->observeAndReduce({R}, {1})->marginalize({C})->normalize();

// Observe R = 1 set I = 0
rcptr<Factor> margC_given_RI = jointICMR->observeAndReduce({R, I}, {1, 0})->marginalize({C})->normalize();
\end{lstlisting}

These results fully correspond to the theory of d-separation. An unobserved descendant of a collider keeps the path blocked, while observing the descendant opens the path and allows statistical influence to flow between the parent nodes.

\section{Question 2}

\subsection*{2.1 (b) Querying the Model}
By using factor operations, the belief that the student will receive a good recommendation letter given that they got a good SAT score and that the course was difficult, or $p(L|D=1,S=1)$, was calculated.

\begin{lstlisting}[language=C++, caption=Calculating recommendation letter probabilities]
// 2.1 b
rcptr<Factor> margL_given_DS = joint->observeAndReduce({D, S}, {1, 1})->marginalize({L})->normalize();
\end{lstlisting}

The output shows that $P(L=0|D=1,S=1) \approx 0.423$ and $P(L=1|D=1,S=1) \approx 0.577$. This means there is a $57.7\%$ chance the student will receive a good recommendation letter. This makes logical sense: even though the course was difficult ($D=1$), the fact that the student achieved a good SAT score ($S=1$) strongly implies higher intelligence ($I=1$). This increased probability of high intelligence makes a good grade (and a good letter) more likely than a bad one.

\subsection*{2.2 (b) Flow of Influence}
In the structure $G \leftarrow I \rightarrow S$, $I$ is a parent of $G$ and $S$. This is a so-called fork structure, and $I$ is called a confounder. We investigate if influence flows between $G$ and $S$ when $I$ is either unobserved or observed.

\begin{itemize}
    \item \textbf{When $I$ is unobserved:} Knowledge of the SAT score $S$ tells us something about the grade $G$. The prior probability of getting the best grade is $P(G=0) = 0.362$. However, when we observe a good SAT score ($S=1$), the probability of getting the best grade jumps to $P(G=0 | S=1) \approx 0.671$. This proves that $S$ and $G$ are dependent.
    
    \item \textbf{When $I$ is observed:} If we already know the student's intelligence (e.g., $I=1$), the probability of getting the best grade is $P(G=0 | I=1) = 0.74$. If we subsequently observe that the student also got a good SAT score ($S=1$), the probability remains exactly $P(G=0 | I=1, S=1) = 0.74$. 
\end{itemize}

\begin{lstlisting}[language=C++, caption=Testing influence in a fork structure]
// Observe S = 1, I unobserved
rcptr<Factor> margG_given_S = joint->observeAndReduce({S}, {1})->marginalize({G})->normalize();

// Observe I = 1, S unobserved
rcptr<Factor> margG_given_I = joint->observeAndReduce({I}, {1})->marginalize({G})->normalize();

// Observe I = 1, S = 1
rcptr<Factor> margG_given_IS = joint->observeAndReduce({I, S}, {1, 1})->marginalize({G})->normalize();
\end{lstlisting}

These conclusions correspond perfectly to the theory of d-separation. In a fork structure, an unobserved confounder allows influence to flow between its children (they are dependent). However, observing the confounder blocks the path, rendering the children conditionally independent of each other.

\section{Question 3: The Casino HMM}

\subsection*{3.1 (d) Querying the Model}
We observed a sequence of rolls: $(Y_1, Y_2, Y_3) = (6, 5, 6)$. \texttt{emdw} was used to calculate the posterior beliefs for the states of the die at each time step.

\begin{lstlisting}[language=C++, caption=Calculating posterior beliefs for the dice states]
// We observe: Rolls 6, 5, 6 (Indices 5, 4, 5)
rcptr<Factor> evidence = jointHMM->observeAndReduce({Y1, Y2, Y3}, {5, 4, 5});

// Calculate P(Z1 | evidence)
rcptr<Factor> postZ1 = evidence->marginalize({Z1})->normalize();

// Calculate P(Z2 | evidence)
rcptr<Factor> postZ2 = evidence->marginalize({Z2})->normalize();

// Calculate P(Z3 | evidence)
rcptr<Factor> postZ3 = evidence->marginalize({Z3})->normalize();
\end{lstlisting}

The output yielded the following probabilities that the die was loaded ($Z=1$) at each time step:
\begin{itemize}
    \item $P(Z_1=1 | Y=(6,5,6)) \approx 0.812$
    \item $P(Z_2=1 | Y=(6,5,6)) \approx 0.771$
    \item $P(Z_3=1 | Y=(6,5,6)) \approx 0.774$
\end{itemize}

These beliefs make sense. A loaded die has a $50\%$ chance of rolling a 6, compared to a fair die's $16.7\%$ chance. Observing two 6s out of three rolls heavily skews the probability toward the die being loaded. 

\subsection*{3.2 (b) Flow of Influence via Chain Nodes}
The structure $Z_1 \rightarrow Z_2 \rightarrow Z_3$ is a chain or cascade, where $Z_2$ acts as a mediator. We investigate whether influence flows between $Z_1$ and $Z_3$ depending on whether $Z_2$ is observed.

\begin{itemize}
    \item \textbf{When $Z_2$ is unobserved:} Knowledge of $Z_3$ tells us something about $Z_1$. The prior is $P(Z_1=1) = 0.5$. However, observing $Z_3=1$ updates this belief to $P(Z_1=1 | Z_3=1) \approx 0.898$. This proves that $Z_1$ and $Z_3$ are dependent.
    
    \item \textbf{When $Z_2$ is observed:} The path is blocked. Observing $Z_2=1$ updates the probability to $P(Z_1=1 | Z_2=1) \approx 0.947$. If we subsequently observe $Z_3=1$, the probability remains exactly $P(Z_1=1 | Z_2=1, Z_3=1) \approx 0.947$. The new information about $Z_3$ does not flow back to $Z_1$.
\end{itemize}

\begin{lstlisting}[language=C++, caption=Testing influence in a chain structure]
// Observe Z3 = 1 (Loaded), Z2 unobserved
rcptr<Factor> margZ1_given_Z3 = jointHMM->observeAndReduce({Z3}, {1})->marginalize({Z1})->normalize();

// Observe Z2 = 1 (Loaded), Z3 unobserved
rcptr<Factor> margZ1_given_Z2 = jointHMM->observeAndReduce({Z2}, {1})->marginalize({Z1})->normalize();

// Observe Z2 = 1 and Z3 = 1
rcptr<Factor> margZ1_given_Z2Z3 = jointHMM->observeAndReduce({Z2, Z3}, {1, 1})->marginalize({Z1})->normalize();
\end{lstlisting}

These conclusions align perfectly with the theory of d-separation. In a chain structure, an unobserved mediator allows influence to flow (dependence), but observing the mediator blocks the path, rendering the start and end nodes conditionally independent.

\end{document}